\documentclass{article}
\usepackage[utf8]{inputenc}
\usepackage{amsmath,amssymb}
\usepackage[most]{tcolorbox}
\usepackage{geometry}
\geometry{margin=2.5cm}

\newcommand{\step}[1]{\par\noindent\hangindent=3.5em\hangafter=1%
  \makebox[3.5em][l]{\textbf{#1.}}}
\newcommand{\steptag}[1]{\hfill{\small\textsf{[#1]}}}

\newtcolorbox{proofbox}[1]{
  colback=gray!5, colframe=gray!60,
  fonttitle=\bfseries, title={#1},
  breakable, enhanced,
  left=4pt, right=4pt, top=4pt, bottom=4pt,
  before skip=10pt, after skip=10pt
}

\begin{document}

\begin{proofbox}{Parameterized approximate-group transport: there exist C\_...}
\step{1} Parameterized approximate-group transport: there exist C\_grp\textasciicircum{}0, C\_pol\textasciicircum{}0 >= 1 and kappa\_pol\textasciicircum{}0 in (0, 1/2] such that, for every def-stage1-polar-witness-data tuple W with C\_grp >= C\_grp\textasciicircum{}0, C\_pol >= C\_pol\textasciicircum{}0, and 0 < kappa\_pol <= kappa\_pol\textasciicircum{}0, for every finite-dimensional exact-unit epsilon\_r-C*-algebra and every delta > 0 satisfying C\_pol*(epsilon\_r + delta) <= kappa\_pol and C\_grp*epsilon\_r < delta - C\_pol*(epsilon\_r*delta + delta\textasciicircum{}2), writing (u\_delta, h\_delta) for the unique inverse of Pi\_delta: calU x B\textasciicircum{}\{calH\}\_delta(J) -> S\_delta := Pi\_delta(calU x B\textasciicircum{}\{calH\}\_delta(J)), Pi\_delta(U, H) = U bold-dot H, the formulas mu(U, V) = u\_delta(U bold-dot V) and sigma(U) = u\_delta(U\textasciicircum{}dagger) define C\textasciicircum{}1 maps on all of calU x calU and calU, respectively, and for every U, V, Z in calU, mu(J, U) = mu(U, J) = U, sigma(J) = J, ||mu(U, V) - U bold-dot V|| <= C\_grp*epsilon\_r, ||sigma(U) - U\textasciicircum{}dagger|| <= C\_grp*epsilon\_r, ||mu(mu(U, V), Z) - mu(U, mu(V, Z))|| <= C\_grp*epsilon\_r, ||mu(sigma(U), U) - J|| <= C\_grp*epsilon\_r, and ||mu(U, sigma(U)) - J|| <= C\_grp*epsilon\_r. \steptag{validated} \\[6pt]
\step{1.1} Provider witnesses and transported constants. Invoke lem-stage1-explicit-group-domain-membership once to fix universal G\_d,P\_d >= 1 and k\_d in (0,1/2]; invoke lem-stage1-explicit-group-closeness once to fix universal G\_c,P\_c >= 1 and k\_c in (0,1/2]; and invoke lem-stage1-polar-retraction once to fix universal P\_r >= 1 and k\_r in (0,1/2]. Define C\_grp\textasciicircum{}0=max\{G\_d,8G\_c,8\}, C\_pol\textasciicircum{}0=max\{P\_d,P\_c,P\_r\}, and kappa\_pol\textasciicircum{}0=min\{k\_d,k\_c,k\_r,1/16\}. Then C\_grp\textasciicircum{}0,C\_pol\textasciicircum{}0 >= 1 and 0<kappa\_pol\textasciicircum{}0<=1/2, and these constants are fixed before any receiving def-stage1-polar-witness-data tuple is quantified. \steptag{validated} \\[4pt]
\step{1.2} Transport of every provider guard. Fix an arbitrary receiving tuple W, finite-dimensional exact-unit epsilon\_r-C*-algebra, and delta as in node 1, using the constants of 1.1; write e=epsilon\_r and q=e*delta+delta\textasciicircum{}2. Since P\_d,P\_c,P\_r <= C\_pol and kappa\_pol <= k\_d,k\_c,k\_r, the receiving inequality C\_pol*(e+delta)<=kappa\_pol gives P\_i*(e+delta)<=k\_i for i=d,c,r. Since G\_d,G\_c<=C\_grp and P\_d,P\_c<=C\_pol, C\_grp*e<delta-C\_pol*q gives G\_i*e<delta-P\_i*q for i=d,c. Also C\_pol>=1, delta>0, and C\_pol*(e+delta)<=kappa\_pol<=1/16 imply 0<=e<1/16. Thus all three fixed provider hypotheses hold for this same algebra and delta. \steptag{validated} \\[4pt]
\step{1.3} Typed conclusions of the three allowed providers. Under 1.2, lem-stage1-explicit-group-domain-membership supplies its unique inverse (u\_d,h\_d) of the displayed Pi\_delta and says every U bold-dot V and U\textasciicircum{}dagger lies in S\_delta and has a right inverse. Under the same hypotheses, lem-stage1-explicit-group-closeness supplies its unique inverse (u\_c,h\_c) and the bounds ||u\_c(U bold-dot V)-U bold-dot V||<=G\_c*e and ||u\_c(U\textasciicircum{}dagger)-U\textasciicircum{}dagger||<=G\_c*e. Finally lem-stage1-polar-retraction supplies a C\textasciicircum{}1 diffeomorphism Pi\_delta from calU x B\_delta\textasciicircum{}\{calH\}(J) onto S\_delta, with unique C\textasciicircum{}1 inverse (u\_r,h\_r), u\_r(U)=U and h\_r(U)=J for U in calU. \steptag{validated} \\[4pt]
\step{1.4} Synchronization by ordinary inverse uniqueness. In all three provider conclusions of 1.3, Pi\_delta has literally the same displayed formula Pi\_delta(U,H)=U bold-dot H, the same domain calU x B\_delta\textasciicircum{}\{calH\}(J), and the same image S\_delta:=Pi\_delta(calU x B\_delta\textasciicircum{}\{calH\}(J)). Hence uniqueness of the inverse of this one bijection gives (u\_d,h\_d)=(u\_c,h\_c)=(u\_r,h\_r). They therefore equal the (u\_delta,h\_delta) named in node 1; in particular every domain, closeness, C\textasciicircum{}1, and retraction assertion in 1.3 concerns that same u\_delta. \steptag{validated} \\[4pt]
\step{1.5} Elementary uniform norm estimates used below. For every T in calU, def-approximate-unitary-space gives T\textasciicircum{}dagger bold-dot T=J and a right inverse R with T bold-dot R=J. The C*-lower bound in def-epsilon-cstar-algebra and ||J||=1 give ||T||<=a:=(1-e)\textasciicircum{}(-1/2). Therefore, for 0<=e<1/16, multiplying a perturbation on either side by a unitary costs at most (1+e)a<2, and the associator of three unitaries has norm at most e*a\textasciicircum{}3<=2e (with equality possible at e=0). Moreover exact unitality and the associator bound give ||R-T\textasciicircum{}dagger||=||(T\textasciicircum{}dagger bold-dot T) bold-dot R-T\textasciicircum{}dagger bold-dot (T bold-dot R)||<=e||T||\textasciicircum{}2||R||, whence ||R||<=||T||(1-e)/(1-2e); consequently ||T bold-dot T\textasciicircum{}dagger-J||<=e(1+e)/((1-e)(1-2e))<=2e (again allowing equality at e=0). These estimates also apply when T is mu(A,B) or sigma(A), since u\_delta takes values in calU. \steptag{validated} \\[4pt]
\step{1.6} Global definition and C\textasciicircum{}1 regularity of the operations. By 1.3-1.4, U bold-dot V and U\textasciicircum{}dagger lie in S\_delta for every U,V in calU, so mu(U,V)=u\_delta(U bold-dot V) and sigma(U)=u\_delta(U\textasciicircum{}dagger) are defined on all of calU x calU and calU and take values in calU. By lem-stage1-polar-retraction, u\_delta is C\textasciicircum{}1 on S\_delta; bilinearity of bold-dot and conjugate-linearity of dagger from def-epsilon-cstar-algebra make the two input maps C\textasciicircum{}1 (as real maps) in finite dimension. Their compositions mu and sigma are therefore C\textasciicircum{}1. \steptag{validated} \\[4pt]
\step{1.7} Left unit law. For U in calU, exact unitality gives J bold-dot U=U, and the retraction identity synchronized in 1.4 gives u\_delta(U)=U. Hence mu(J,U)=u\_delta(J bold-dot U)=U. \steptag{validated} \\[4pt]
\step{1.8} Right unit law. For U in calU, exact unitality gives U bold-dot J=U, and the retraction identity synchronized in 1.4 gives u\_delta(U)=U. Hence mu(U,J)=u\_delta(U bold-dot J)=U. \steptag{validated} \\[4pt]
\step{1.9} Unit inversion law. Exact-unit compatibility with the involution in def-epsilon-cstar-algebra gives J\textasciicircum{}dagger=J, while 1.4 gives u\_delta(J)=J. Thus sigma(J)=u\_delta(J\textasciicircum{}dagger)=J. \steptag{validated} \\[4pt]
\step{1.10} Product closeness. The synchronized closeness provider in 1.3-1.4 gives ||mu(U,V)-U bold-dot V||=||u\_delta(U bold-dot V)-U bold-dot V||<=G\_c*e. Since C\_grp>=C\_grp\textasciicircum{}0>=8G\_c>=G\_c, this is at most C\_grp*epsilon\_r. \steptag{validated} \\[4pt]
\step{1.11} Inverse closeness. The synchronized closeness provider in 1.3-1.4 gives ||sigma(U)-U\textasciicircum{}dagger||=||u\_delta(U\textasciicircum{}dagger)-U\textasciicircum{}dagger||<=G\_c*e. Since C\_grp>=C\_grp\textasciicircum{}0>=8G\_c>=G\_c, this is at most C\_grp*epsilon\_r. \steptag{validated} \\[4pt]
\step{1.12} Approximate associativity of mu. Put A=mu(mu(U,V),Z) and B=mu(U,mu(V,Z)). By 1.10 at the provider strength G\_c, the two outer retraction errors contribute G\_c*e each. Replacing mu(U,V) by U bold-dot V and mu(V,Z) by V bold-dot Z inside one multiplication contributes less than 2G\_c*e on each side by 1.5, and the intervening algebra associator contributes less than 2e. Hence ||A-B||<=(6G\_c+2)e. From C\_grp>=8G\_c and C\_grp>=8, one has 6G\_c+2<=C\_grp, proving ||mu(mu(U,V),Z)-mu(U,mu(V,Z))||<=C\_grp*epsilon\_r. \steptag{validated} \\[4pt]
\step{1.13} Left approximate inverse law. Since sigma(U) is in calU, product closeness gives ||mu(sigma(U),U)-sigma(U) bold-dot U||<=G\_c*e. By inverse closeness and the non-strict multiplication estimate in the amended node 1.5, ||sigma(U) bold-dot U-U\textasciicircum{}dagger bold-dot U||<=(1+e)*||sigma(U)-U\textasciicircum{}dagger||*||U||<=2G\_c*e; this remains valid when e=0. Also U\textasciicircum{}dagger bold-dot U=J by def-approximate-unitary-space. Hence the triangle inequality gives ||mu(sigma(U),U)-J||<=3G\_c*e<=C\_grp*e, because C\_grp>=C\_grp\textasciicircum{}0>=8G\_c. \steptag{validated} \\[4pt]
\step{1.14} Right approximate inverse law. Product closeness gives ||mu(U,sigma(U))-U bold-dot sigma(U)||<=G\_c*e. By inverse closeness and the multiplication estimate of 1.5, ||U bold-dot sigma(U)-U bold-dot U\textasciicircum{}dagger||<=(1+e)*||U||*||sigma(U)-U\textasciicircum{}dagger||<=2G\_c*e; this remains non-strict when e=0. The defect estimate derived in 1.5 gives ||U bold-dot U\textasciicircum{}dagger-J||<=e*(1+e)/((1-e)*(1-2e))<=2e for every 0<=e<1/16, including e=0. Hence the triangle inequality yields ||mu(U,sigma(U))-J||<=(3G\_c+2)*e<=C\_grp*e: indeed C\_grp>=8G\_c gives 3G\_c<=(3/8)C\_grp and C\_grp>=8 gives 2<=C\_grp/4, so 3G\_c+2<=(5/8)C\_grp<=C\_grp. \steptag{validated} \\[4pt]
\step{1.15} Assembly. Nodes 1.1-1.14 provide universal constants of the required ranges, transport all hypotheses to the three allowed providers, identify their inverse components with the single (u\_delta,h\_delta) bound in the contract, establish global C\textasciicircum{}1 maps, and prove every displayed identity and estimate for arbitrary W, algebra, delta, and U,V,Z. Universal generalization and the definitions of mu and sigma therefore yield node 1 exactly. \steptag{validated} \\[4pt]
\end{proofbox}

\end{document}
